\section{Conclusions and Future Work} %no future work so far (not strictly needed) %  and Future Work} \todo{[Claus,Jens]}
\label{sec:conclusion}
%\vspace{-5mm}
In this article, we presented work towards a unified model for RDB2RDF mappings and the
lightweight mapping language \emph{SML} that reuses familiar elements of SPARQL
and SQL in order to lower the learning curve and ease the manual writing and maintenance of view definitions.
An extensive public survey confirmed that this is the case.
%\todo{R3: would need to be supported with a usability study}
We provided an in-depth comparison of how SML relates to the R2RML standard, and
detailed how the former can be automatically converted to the latter.
%A comprehensive formalization of the RDB-RDF rewriting process based the
%definition of RDF views of a relational database and bindings, which bridge
%between the relational and SPARQL algebra.

SML has been successfully deployed in several scenarios:
We created SML mappings for the BSBM and SP2 benchmarks, two
popular SPARQL benchmarks that are often used for evaluating RDB2RDF
mappers.
Most prominently, we created SML mappings for transforming the
OpenStreetMap (OSM) database to RDF. These efforts are carried out
as part of the LinkedGeoData\footnote{\url{http://linkedgeodata.org/}} (LGD) project,
where we give access to more than 25 billion OSM RDF triples created through
SPARQL-to-SQL rewriting over about 3 billion relational rows via more than 40
SML view definitions.

%LGD also includes mappings for the GADM dataset
Furthermore, SML mappings have been created for two large scale
linguistic resources: One is the mapping of the Wortschatz
database\footnote{\url{http://www.wortschatz.uni-leipzig.de/}}, which contains
statistics, such as frequency and co-occurrences, about words in more
than 240 languages.
The other resource is PanLex\footnote{\url{http://ld.panlex.org}}, which is
a database holding translations of about 19 million expression extracted from
over 2.000 sources. Links to the corresponding SML mappings are published
together with our other SML related
resources\footnote{\url{http://sml.aksw.org}}.

In general, we believe that mapping relational structures to RDF will stay a highly important topic in research and practice to provide an unobtrusive transition towards the use of semantic technologies.
Providing engineers an intuitive yet powerful language is a crucial step to ease this transition.
Future work will continue on extending the formalizations as well as sorting out details based on community feedb,
such as whether an explicit \texttt{FROM QUERY} syntax for specifying SQL queries is preferred over
the current approach where this is implied by the use of triple quotes.


% SPARKLIFY

%\section{Conclusions and Future Work}
%\label{sec:conclusion}
Querying RDF data becomes challenging when the size of the data increases.
%Existing systems mostly contain a very moderately information of the data while transforming them to a dedicated storage model (e.g. vertical partitioning).
%contain a very moderately information of the data
Existing Spark-based SPARQL systems mostly
do not retain all RDF term information consistently while transforming them to a dedicated storage model such as using vertical partitioning.
Often, this process is both data and computing intensive and raises the need for a scalable, efficient and comprehensive query engine which can handle large scale RDF datasets.

In this paper, we propose \emph{Sparklify}: a scalable software component for efficient evaluation of SPARQL queries over distributed RDF datasets.
It uses Sparqify as a SPARQL-to-SQL rewriter for translating SPARQL queries into Spark executable code.
By doing so, it leverages the advantages of the Spark framework.
SANSA features methods to execute SPARQL queries directly as part of Spark workflows instead of writing the code corresponding to those queries (sorting, filtering, etc.).
It also provides a command-line interface and a W3C standard compliant SPARQL endpoint for externally querying data that has been loaded using the SANSA framework.
We have shown empirically that our approach can scale horizontally and perform well w.r.t to the state-of-the-art approaches.

With this work, we showed that the application of OBDA tooling to Big Data frameworks achieves promising results in terms of scalability.
We present a working prototype implementation that can serve as a baseline for further research.
Our next steps include evaluating other tools, such as Ontop~\cite{Calvanese2017OntopAS}, and analyze how their performance in the Big Data setting can be improved further.
For example, we intend to investigate how OBDA tools can be combined with dictionary encoding of RDF terms as integers and evaluate the effects.

%\section{Conclusions and Future Work}
%\label{sec:conclusions}
% Sansa-RML
In this work we showed that with the conversion of RML to SPARQL we can leverage suitably enhanced SPARQL engines for the task of knowledge graph construction.
%Just like YARRRML can be converted to RML,  a common ground for query and workload optimization.
We further showed that by transforming CONSTRUCT queries to their ``lateral'' form it is now finally possible to ``merge'' CONSTRUCT queries and remove duplicates which has direct applications in knowledge graph construction.
%Furthermore, using a simple join analysis we can optimize JOINs on Spark by transforming expensive hash joins into cheaper broadcast joins.
Using query workload analysis we can push down DISTINCT operations such that this expensive operation can be computed on smaller RDF graphs.
We showed that the same query workload can be executed on different engines yielding the same result sets however with significantly different performance characteristics. By leveraging a Big Data framework this approach can outperform state of the art approaches.
%In particular, to the best of our knowledge, our Big Data-based approach is the first one to scale to the GTFS1000 benchmark.
We emphasize that as part of this work we contributed the SERVICE extension plugin system as well as minor performance improvements
%as an initial LATERAL implementation
to Apache Jena.
%Our contribution was picked up in Oxigraph and sparked standardization efforts for SPARQL 1.2.
One direction of future work is to optimize the generated SPARQL algebra further as to minimize the amount of data that has to be processed in DISTINCT and ORDER BY operations. Also, as shown in the evaluation, the improved overall performance comes at the cost of significant higher resource usage for which we plan in-depth investigation of the reasons and possible mitigation approaches such as using custom Spark operator implementations.
Furthermore, we identify the need for the standardization of SPARQL for heterogeneous data as this would not only make it possible to transform RML to SPARQL in a truly interoperable way, but also provide a common ground for query and query workload optimization.

